\chapter{总结与展望}

\textbf{【第七章撰写总纲】} 本章对全文进行收束：7.1 按“问题—方法—结论”梳理研究贡献并与摘要、各章小结呼应；7.2 客观指出当前工作的局限（数据、规模、场景等），避免过度宣称；7.3 基于局限与领域趋势提出 2–4 个可操作的未来方向，与第六、七章可自然衔接（如自动模型选择、多目标优化等）。篇幅上建议 7.1 约 2–3 页，7.2、7.3 各约 1 页。

---

\section{研究工作总结}

\textbf{【撰写大纲】}

1. \textbf{研究背景与问题回顾}
\begin{itemize}
\item 用 1 段概括：多元长程时间序列预测在能源、交通、金融等领域的意义；现有方法在“时频兼顾、LLM 融合、计算效率”等方面的不足；本文聚焦的核心问题——如何在可控开销下融合大语言模型与时频信息、并兼顾不同评价指标下的表现。
\end{itemize}

2. \textbf{主要工作与贡献（按章归纳）}
\begin{itemize}
\item \textbf{统一框架（第三章）}：提出基于 Qwen3 与门控注意力的统一改进框架；冻结 LLM + 轻量适配器实现时序与语义对齐；门控注意力池化（编码层 SDPA 门控 + 多头 CMA 门控融合）抑制冗余、提升稳定性；基线实验验证框架有效性与 Qwen3 相对 GPT-2、门控机制的收益。
\item \textbf{小波包时频域增强（第四章）}：在统一框架下引入小波包多尺度分解；给出分解策略、子带与 LLM 映射、多尺度融合与预测头的完整设计；实验表明在 \textbf{MSE} 及极端波动场景下优于基线，消融支持当前分解层数配置。
\item \textbf{频域前馈网络增强（第五章）}：在统一框架下引入 FreTS 频域前馈网络；给出复数域线性、稀疏化与频域-时域交互设计；实验表明该全局频域增强方法在综合预测性能、长期趋势稳定性与平均误差控制方面更具优势，消融实验进一步支持当前配置的合理性。
\item \textbf{机理与模型选择（第六章）}：对局部多尺度频域增强路线与全局频域增强路线进行综合对比；从误差特性与频域设计角度解释两类方法在不同数据特性和评价指标下的表现差异；给出计算复杂度、收敛性对比及基于误差敏感度的场景分类与模型选择建议。
\end{itemize}

3. \textbf{总体结论}
\begin{itemize}
\item 用 1 段总结：本文通过统一框架将 LLM、时域、两种频域路线与门控机制有机结合，在多个基准上取得较优的 MSE/MAE 表现，并揭示了不同频域设计对指标分化的影响，为实际应用中按业务需求与数据特性选择预测模型提供了理论依据与实验支撑；与摘要、摘要结论保持一致表述。
\end{itemize}

---

\section{研究局限性分析}

\textbf{【撰写大纲】}

1. \textbf{数据与场景局限}
\begin{itemize}
\item 实验主要基于公开基准数据集（ETT、Weather、Exchange、ILI 等），\textbf{领域与规模有限}；在特定行业（如电力、医疗、金融实盘）或超长序列、超高维变量场景下的泛化能力尚未充分验证。
\item 若存在\textbf{缺失值、非均匀采样、多粒度时间}等复杂情况，当前预处理与模型是否仍适用需进一步研究。
\end{itemize}

2. \textbf{模型与计算局限}
\begin{itemize}
\item \textbf{LLM 规模}：本文采用 Qwen3-0.6B 作为外部语义编码器，更大或更小规模 LLM 对性能与效率之间的权衡关系尚未系统探索；\textbf{冻结策略}在少样本或领域迁移场景下是否仍为最优仍可进一步讨论。
\item \textbf{频域路线}：WPD 与 FreTS 为两条独立分支，\textbf{联合或动态切换}（如根据输入自动选择频域路径）尚未在本文实现；二者在统一损失下的\textbf{多目标或自适应加权}也未深入展开。
\item \textbf{计算资源}：在资源受限设备或实时场景下，推理延迟与显存占用是否满足需求，依赖具体部署与优化，本文未做系统评估。
\end{itemize}

3. \textbf{评价与业务局限}
\begin{itemize}
\item 评价以 \textbf{MSE、MAE} 为主，其他指标（如分位数损失、峰值误差、下游决策损失）在不同业务中的重要性未充分讨论；\textbf{模型选择策略}（第六章）主要基于误差敏感度与数据特性，尚未与具体业务 KPI 或成本函数做定量挂钩。
\end{itemize}

4. \textbf{表述注意}
\begin{itemize}
\item 用客观、简洁的语言陈述局限，避免自我贬低；每点 1–3 句即可，为 7.3 未来方向做自然过渡。
\end{itemize}

---

\section{未来研究方向}

\textbf{【撰写大纲】}

1. \textbf{自适应频域与模型选择}
\begin{itemize}
\item 研究\textbf{输入依赖的频域路径选择}：根据序列特性（如方差、频谱熵、突变检测）自动在 WPD 与 FreTS（或更多频域模块）间切换或加权，减少人工选择成本；与 6.4 模型选择策略结合，探索轻量级“选择器”网络或规则。
\item \textbf{多目标与 Pareto 优化}：在训练中显式考虑 MSE 与 MAE（或更多指标）的 Pareto 前沿，生成一组模型供部署时按业务权重选择；或设计单一模型的多目标损失，平衡不同误差类型。
\end{itemize}

2. \textbf{更大规模 LLM 与少样本/领域迁移}
\begin{itemize}
\item 在\textbf{更大参数量 LLM}（如 7B、14B）或\textbf{领域预训练/微调}下，评估统一框架的扩展性与收益；探索\textbf{少样本或零样本}场景下适配器与门控机制的作用，以及\textbf{跨领域迁移}（如从电力到交通）的泛化表现。
\end{itemize}

3. \textbf{复杂数据与下游任务}
\begin{itemize}
\item 针对\textbf{缺失值、非等间隔、多粒度时间}的建模与预处理；将预测模块与\textbf{下游决策任务}（如调度、告警、资源配置）联合优化，以业务损失或约束为导向设计损失与评估指标，与 7.2 的“业务 KPI”局限呼应。
\end{itemize}

4. \textbf{效率与部署}
\begin{itemize}
\item \textbf{模型压缩与蒸馏}：在保持精度的前提下压缩统一框架（含 LLM 或仅时序/频域部分），便于边缘或实时部署；\textbf{推理加速}（如算子融合、量化、早停）在长序列与多变量下的应用。
\end{itemize}

5. \textbf{撰写建议}
\begin{itemize}
\item 从中选取 2–4 点重点展开，每点 1 段；与 7.2 的局限一一对应更佳（如“数据局限→复杂数据与下游任务”“模型选择局限→自适应频域与模型选择”）；结尾可用 1–2 句总括“上述方向将进一步完善本文框架的适用性与实用性”。
\end{itemize}

---
